% Divisors and the Degree-Zero Theorem  (source §5; added: the divisor group
% Div(X), degree homomorphism, linear equivalence and Pic^0, and the Abel-sum
% map Div^0(T) -> T with Pic^0(T) =~ T)
\chapter{Divisors and the Degree-Zero Theorem}
\label{ch:divisors}

\Cref{ch:elliptic_liouville} controlled the zeros and poles of a single elliptic
function. Divisors are the bookkeeping that turns those zero/pole data into an
algebraic object one can add and compare. On a compact Riemann surface the
divisor of a function has degree zero; on the torus it additionally has Abel sum
zero. Packaging these two constraints gives the \emph{Picard group}
\(\Pic^0(T)\) and a map \(\Pic^0(T) \to T\) that is the structural shadow of the
group law: when we later put a group law on the cubic, it will be exactly this
map.

\section{Orders and the Divisor of a Function}
\label{sec:div-orders}

Let \(X\) be a Riemann surface and \(f\) a nonzero meromorphic function. In a
local coordinate \(t\) centered at a point \(p\) we may write
\[
  f(t) = t^{m} u(t), \qquad u \text{ holomorphic}, \quad u(0) \ne 0 .
\]
The integer \(m\) is the \emph{order} of \(f\) at \(p\), written
\(\ord_p(f) = m\):
\[
  m > 0 \ \text{(zero of order } m),\qquad
  m = 0 \ \text{(finite, nonzero)},\qquad
  m < 0 \ \text{(pole of order } {-m}).
\]

\begin{remark}
\label{rem:div-ord-welldef}
The order does not depend on the chart. If \(s\) is another coordinate at
\(p\), then \(t = s\,h(s)\) with \(h\) holomorphic and \(h(0) \ne 0\), so
\(f = s^{m}\bigl(h(s)^{m} u(t(s))\bigr)\) and the new unit
\(h^{m}\,(u\circ t)\) is again nonvanishing at \(s=0\); the exponent \(m\) is
unchanged. Thus \(\ord_p(f)\) is intrinsic.
\end{remark}

\begin{definition}
\label{def:div-divisor-of-f}
The \emph{divisor} of a nonzero meromorphic function \(f\) is the formal sum
\[
  \divisor(f) = \sum_{p \in X} \ord_p(f)\, p .
\]
On a compact \(X\) only finitely many terms are nonzero, so this is a finite
sum.
\end{definition}

\section{The Divisor Group and Degree}
\label{sec:div-group}

\begin{definition}
\label{def:div-group}
A \emph{divisor} on \(X\) is a finite formal integer combination of points,
\[
  D = \sum_{p \in X} n_p\, p, \qquad n_p \in \Z, \ \ n_p = 0 \text{ for all but finitely many } p .
\]
Under coefficientwise addition the divisors form the free abelian group
\(\Div(X)\) on the points of \(X\). The \emph{degree} is the homomorphism
\[
  \degdiv : \Div(X) \to \Z, \qquad \degdiv\Bigl( \sum_p n_p\, p \Bigr) = \sum_p n_p .
\]
A divisor of the form \(\divisor(f)\) is called \emph{principal}; we write
\(\Prin(X)\) for the set of principal divisors.
\end{definition}

\begin{proposition}
\label{prop:div-prin-subgroup}
The principal divisors form a subgroup of \(\Div(X)\): for nonzero meromorphic
\(f, g\),
\[
  \divisor(fg) = \divisor(f) + \divisor(g),
  \qquad
  \divisor(1/f) = -\divisor(f),
\]
and \(\divisor(c) = 0\) for a nonzero constant \(c\).
\end{proposition}

\begin{proof}
Orders add under multiplication: \(\ord_p(fg) = \ord_p(f) + \ord_p(g)\), directly
from \(f = t^{m}u\), \(g = t^{n}v\) giving \(fg = t^{m+n}uv\). Summing over \(p\)
gives the first identity; the others are the cases \(g = 1/f\) and \(f = c\).
Hence \(\Prin(X)\) is closed under addition and negation, a subgroup.
\end{proof}

\section{Principal Divisors Have Degree Zero}
\label{sec:div-degzero}

The single most important property of principal divisors on a compact surface is
that they are invisible to the degree.

\begin{theorem}[Degree-zero theorem]
\label{thm:div-degzero}
If \(X\) is compact and \(f\) is a nonzero meromorphic function on \(X\), then
\[
  \degdiv \divisor(f) = 0 .
\]
\end{theorem}

\begin{proof}
Consider the meromorphic differential \(df/f\). Near \(p\), writing
\(f = t^{m} u\) with \(u(0) \ne 0\),
\[
  \frac{df}{f} = \frac{d(t^{m}u)}{t^{m}u}
  = m\,\frac{dt}{t} + \frac{du}{u},
\]
and the second term is holomorphic at \(p\). Hence \(df/f\) has at worst a simple
pole at \(p\), with
\[
  \Res_p\!\left( \frac{df}{f} \right) = m = \ord_p(f) .
\]
By the residue theorem on the compact Riemann surface \(X\) the sum of all
residues of the global differential \(df/f\) vanishes, so
\[
  0 = \sum_{p} \Res_p\!\left( \frac{df}{f} \right) = \sum_p \ord_p(f)
  = \degdiv \divisor(f) . \qedhere
\]
\end{proof}

\begin{keyidea}{What degree zero means}
A meromorphic function on a compact Riemann surface has as many zeros as poles,
counted with multiplicity --- the divisor form of the argument principle. On the
torus this is \cref{thm:ell-valence} with \(N = \deg(\text{pole divisor})\): the
two theorems are the same fact seen from the two sides of the
\(T \leftrightarrow \C\) dictionary.
\end{keyidea}

In particular \(\Prin(X) \subseteq \Div^0(X)\), where
\[
  \Div^0(X) = \ker(\degdiv) = \{ D : \degdiv D = 0 \}
\]
is the group of degree-zero divisors.

\section{Linear Equivalence and the Picard Group}
\label{sec:div-picard}

\begin{definition}
\label{def:div-lineq}
Two divisors \(D, D'\) are \emph{linearly equivalent}, written \(D \sim D'\), if
\(D - D'\) is principal. The quotient group
\[
  \Pic(X) = \Div(X)/\Prin(X)
\]
is the \emph{Picard group}, and the class of \(D\) is denoted \([D]\).
\end{definition}

Because principal divisors have degree zero (\cref{thm:div-degzero}), the degree
descends to a homomorphism \(\degdiv : \Pic(X) \to \Z\), and its kernel
\[
  \Pic^0(X) = \Div^0(X)/\Prin(X)
\]
is the \emph{degree-zero Picard group}. The size of \(\Pic^0\) measures how far
divisors of degree zero are from being principal --- equivalently, how many
meromorphic functions the surface admits. On \(\CP^1\) every degree-zero divisor
is principal and \(\Pic^0 = 0\); on the torus, as we now show, \(\Pic^0\) is as
large as \(T\) itself.

\section{The Abel--Jacobi Sum Map}
\label{sec:div-abel-jacobi}

Specialize to \(X = T = \C/\Lambda\). Here points of \(T\) are themselves group
elements, so a degree-zero divisor can be \emph{summed}.

\begin{definition}
\label{def:div-sigma}
The \emph{Abel--Jacobi sum map} is
\[
  \smap : \Div^0(T) \to T,
  \qquad
  \smap\Bigl( \sum_i n_i\,[z_i] \Bigr) = \sum_i n_i\,[z_i] \in T,
\]
where the right-hand sum is taken in the group \(T\).
\end{definition}

It is a group homomorphism, immediately from the definition. The content of the
previous chapter is exactly that \(\smap\) annihilates principal divisors.

\begin{proposition}
\label{prop:div-sigma-kills-prin}
For every nonzero elliptic function \(f\), \(\smap(\divisor f) = 0\). Hence
\(\smap\) descends to a homomorphism
\[
  \bar\smap : \Pic^0(T) \to T .
\]
\end{proposition}

\begin{proof}
Write \(\divisor(f) = \sum_k [z_k] - \sum_k [p_k]\) with the zeros and poles
listed by multiplicity; it has degree zero by \cref{thm:div-degzero} (or
\cref{thm:ell-valence}), so it lies in \(\Div^0(T)\). Then
\(\smap(\divisor f) = \sum_k [z_k] - \sum_k [p_k]\), which is \(0\) in \(T\) by
the Abel relation \cref{thm:ell-abel}. Since \(\smap\) kills the subgroup
\(\Prin(T)\), it factors through \(\Pic^0(T) = \Div^0(T)/\Prin(T)\).
\end{proof}

\begin{theorem}[Abel--Jacobi for the torus]
\label{thm:div-pic-iso}
The map \(\bar\smap : \Pic^0(T) \to T\) is an isomorphism of groups. In
particular every complex torus is its own \emph{Jacobian}:
\(\Pic^0(T) \cong T\).
\end{theorem}

\begin{proof}[Proof of surjectivity, and the standing of injectivity]
\emph{Surjectivity.} Given any \([a] \in T\), the divisor \([a] - [0]\) has degree
zero and \(\smap([a]-[0]) = [a] - [0] = [a]\). So \(\bar\smap\) hits every
element of \(T\).

\emph{Well-definedness and homomorphism} are \cref{prop:div-sigma-kills-prin}.

\emph{Injectivity} is the converse of Abel's relation: if \(D \in \Div^0(T)\)
has \(\smap(D) = 0\), then \(D\) is principal, i.e.\ \(D = \divisor(f)\) for some
elliptic \(f\). This requires \emph{constructing} an elliptic function with
prescribed zeros and poles, which we do not yet have the tools for. It follows
once the Weierstrass machinery is in place: \cref{ch:riemann_roch} produces
enough functions, and the explicit construction is carried out through
\crefrange{ch:weierstrass_wp}{ch:transport}. We therefore record injectivity
here and complete it there.
\end{proof}

\begin{keyidea}{Why this is the whole story}
\Cref{thm:div-pic-iso} says the degree-zero divisor classes on the torus
\emph{are} the points of the torus, via summation. Two divisors are linearly
equivalent precisely when their points have the same Abel sum. When we transport
the picture to the cubic \(E\) in \crefrange{ch:lines_divisors}{ch:transport},
the group law on \(E\) will be defined by linear equivalence of divisors --- and
by this theorem that is the same as addition on \(T\). The chord--tangent law is
the composite
\[
  E \;\cong\; T \;\xleftarrow{\ \bar\smap\ }\; \Pic^0(T) ,
\]
read as a single group.
\end{keyidea}

\section{Summary}
\label{sec:div-summary}

Divisors on \(X\) form the free abelian group \(\Div(X)\), graded by degree, with
the principal divisors \(\divisor(f)\) as a subgroup. On a compact surface every
principal divisor has degree zero (\cref{thm:div-degzero}), so degree descends to
the Picard group \(\Pic(X) = \Div(X)/\Prin(X)\) and cuts out \(\Pic^0(X)\). On
the torus the Abel--Jacobi map \(\smap : \Div^0(T) \to T\) sums a divisor's
points; it kills principal divisors by the Abel relation
(\cref{prop:div-sigma-kills-prin}), and the induced
\(\bar\smap : \Pic^0(T) \to T\) is an isomorphism (\cref{thm:div-pic-iso}, with
injectivity completed later). This identification \(\Pic^0(T) \cong T\) is the
algebraic form of the group law that the rest of the book transports to the
cubic.

\section*{Exercises}
\addcontentsline{toc}{section}{Exercises}

\begin{exercise}
\label{ex:div-hom}
Show that \(f \mapsto \divisor(f)\) is a homomorphism from the multiplicative
group of nonzero meromorphic functions to \(\Div(X)\), and that on a compact
\(X\) its kernel is exactly the nonzero constants. (Use
\cref{thm:div-degzero} and \cref{thm:ell-const} for the torus case.)
\end{exercise}

\begin{exercise}
\label{ex:div-ord-welldef}
Prove carefully that \(\ord_p(f)\) is independent of the local coordinate
(\cref{rem:div-ord-welldef}), and deduce that \(\ord_p(f+g) \ge
\min(\ord_p f, \ord_p g)\) with equality when the two orders differ.
\end{exercise}

\begin{exercise}
\label{ex:div-pic-split}
Using the base point \([0] \in T\), show that
\[
  \Div(T) \xrightarrow{\ \sim\ } \Z \oplus \Div^0(T),
  \qquad D \mapsto \bigl( \degdiv D,\ D - (\degdiv D)[0] \bigr),
\]
and deduce a splitting \(\Pic(T) \cong \Z \oplus \Pic^0(T)\). Combined with
\cref{thm:div-pic-iso}, conclude \(\Pic(T) \cong \Z \oplus T\).
\end{exercise}

\begin{exercise}
\label{ex:div-distinct-points}
Let \([a], [b] \in T\). Using the Abel relation \cref{thm:ell-abel}, prove that
the divisor \([a] - [b]\) is principal \emph{if and only if} \([a] = [b]\).
Contrast this with \(\CP^1\), where \emph{any} two points satisfy
\([a] - [b] = \divisor\!\bigl(\tfrac{z-a}{z-b}\bigr)\). Why does this recover
\cref{cor:ell-order2}?
\end{exercise}

\begin{exercise}
\label{ex:div-sigma-compute}
Take \(\Lambda = \Z + \Z\tau\) and the divisor
\(D = [\tfrac12] + [\tfrac{\tau}{2}] + [\tfrac{1+\tau}{2}] - 3\,[0]\). Compute
\(\degdiv D\) and \(\smap(D) \in T\). Is \(D\) a candidate for a principal
divisor? (Recall the two-torsion points from \cref{ex:tg-two-torsion}.)
\end{exercise}
